\section{The concrete bound}
\label{sec:concrete-bound}

The local certificate uses the exact parameters
\[
 \begin{aligned}
 p&=\frac{499999999}{10^9},&
 \mu&=\frac{405108059767}{10^{12}},&
 w&=\frac{104160447497}{250000000000},\\
 \beta&=\frac{9999999}{10^7},&
 x&=\frac{337503}{10^6},&
 \xi&=x-\frac4{10^{11}}.
 \end{aligned}
\]
\leannote{\leanfile{S7_ParameterChoice/L1_Parameters.lean\#L19}{L1_Parameters.lean}}
The selected terminal source is
\[
 A_s=\frac{1086180879838687}{10^{15}},\qquad
 B_s=\frac{11640399543290787172985714354555752537}
 {36367644171000000000000000000000000000}.
\]
\leannote{\texttt{sourceA} and \texttt{sourceB},
\leanfile{S7_ParameterChoice/L1_Parameters.lean\#L27}{L1_Parameters.lean}}
All remaining terminal parameters, knots, slopes, and profile heights
are specified exactly in the linked certificate files.

\subsection{Certifying the local envelope}

Set $a=2681/10000$. Rational interval bounds give
\[
 \mathcal F(a)>\frac{1769}{10^8}>0,
 \qquad \mathcal F(\mu)<-\frac{47763}{10^6}<0.
\]
\leannote{\texttt{entropy\_gap\_at\_cutoff} and \texttt{entropy\_gap\_at\_endpoint},
\leanfile{S7_ParameterChoice/L2_Control.lean\#L100}{L2_Control.lean}}
Lemma~\ref{lem:feasibility-cutoff} therefore gives
$\mu\in F_p\subset(a,\mu]$. The certificate must cover this whole
interval, rather than just a collection of sampled fractions.

Write $H(q)=qM_p(q)+wq$. Its global upper supports follow from
concavity on the positive-budget domain
$D=\{q\in(0,1):p e^{\beta A_p(q)}<U(q)\}$.
Here is the reason for that concavity. Set
$G_0(q)=q\log p+\beta q(A_p(q)-B_q)$,
$f(z)=pz+(1-p)\log(1-e^z)$, and
$\bar f(z)=f(\max\{z,\log p\})$. Then
\[
 H(q)=qB_q+wq+
 \frac q\beta\left[\bar f\left(\frac{G_0(q)}q\right)-f(\log p)\right].
\]
\leannote{\texttt{certM\_entropy\_identity},
\leanfile{S4_ScalarOptimization/L13_EnvelopeConcavity.lean\#L102}{L13_EnvelopeConcavity.lean}}
The function $G_0$ is convex: its nonlinear part is $\beta w$ times
binary entropy. Thus $D=\{G_0<0\}\cap(0,1)$ is convex.
On $(-\infty,0)$, $\bar f$ is concave and nonincreasing; this follows
from $f'(z)=(p-e^z)/(1-e^z)$ and $f''(z)<0$.
The perspective $q\bar f(G_0(q)/q)$ is consequently concave:
normalize the convex-combination weights by their weighted mean,
then use monotonicity and concavity of $\bar f$.
The remaining term $qB_q=wq\log(\mu/q)$ is concave as well.
The explicit support slopes are in
\leanfile{S4_ScalarOptimization/L13_EnvelopeConcavity.lean\#L351}{L13_EnvelopeConcavity.lean}.

Divide $[a,\mu]$ into $192$ equal cells, with endpoints
$q_i=a+i(\mu-a)/192$ and midpoints $m_i=(q_i+q_{i+1})/2$.
Checked endpoint budget signs place all of $[a,\mu]$ in $D$.
At each midpoint take the explicit upper support $T_i$ of $H$.
The certificate supplies rational heights $h_0,\ldots,h_{192}$ with
\[
 H(q)\le T_i(q)\quad(q\in D),\qquad
 T_i(q_i)\le h_i,\qquad T_i(q_{i+1})\le h_{i+1}.
\]
\leannote{\texttt{cell\_inputs},
\leanfile{S7_ParameterChoice/Control/Envelope.lean\#L97}{Envelope.lean},
and the heights in \leanfile{S7_ParameterChoice/Control/Table.lean\#L26}{Table.lean}}
Since $T_i$ is affine, its endpoint bounds extend throughout the cell.
Adding the affine function $wqv$ then bounds $H(q)+wqv$ by the maximum
of the two endpoint values. This proves coverage for every feasible
$q$ and every real $v$.

Put $a_i=wq_i$, $\Delta_i=a_{i+1}-a_i>0$, and
$r_i=(h_i-h_{i+1})/\Delta_i$. The logarithmic hinge
\[
 E(u)=a_0\log u+h_0+
 \sum_{i=0}^{191}\Delta_i(\log u-r_i)_+
 \qquad(0<u\le1)
\]
\leannote{\texttt{affineHingeMajorant\_dominates},
\leanfile{S4_ScalarOptimization/L9_AffineEnvelope.lean\#L130}{L9_AffineEnvelope.lean}}
dominates every line $a_i\log u+h_i$: telescope the successive line
differences and replace each difference by its positive part.
It is continuous, nondecreasing, and integrable. Integrating a single
hinge gives $\int_0^1(\log u-r)_+\,du\le e^r-r-1$, with equality
for $r\le0$. Telescoping again, the checked exponential sum yields
\begin{equation}\label{eq:certified-integral}
 \int_0^1E(u)\,du
 \le h_{192}-a_{192}+\sum_{i=0}^{191}\Delta_i e^{r_i}
 <-\frac{15}{10^8}.
\end{equation}
\leannote{\texttt{integral\_sum\_lt},
\leanfile{S7_ParameterChoice/Control/Integral.lean\#L102}{Integral.lean};
the integration law is in \leanfile{S4_ScalarOptimization/L9_AffineEnvelope.lean\#L183}{L9_AffineEnvelope.lean}}
This proves the negative-envelope premise of
Proposition~\ref{prop:local-alternative}.
The interval calculations include proofs that the rational expressions
represent the stated real functions and that every logarithm and inverse
is evaluated in its valid domain. The complete assembly is
\leanfile{S7_ParameterChoice/Control/Envelope.lean\#L152}{Envelope.lean}.

\subsection{The joint family and the root}

We reassemble the inherited profile data of~\cite{ZBS} into a supporting
family of $1469$ direct profiles, including five custom profiles,
$947$ combined profiles, and $2716$ source lines.
Their certificates cover every profile domain and every positive source
ratio. The source witnesses have uniform reserve
$1/(2\cdot10^{12})$; the selected source $(A_s,B_s)$ has the stronger
reserve $10^{-9}$, using $81$ profile-pair runs and two baseline tails.
The baseline is Lu--Wang's uniform estimate~\cite[Lemma~B.1]{LW}.
The family and coverage proofs are recorded in
\href{\leanrepo/S7_ParameterChoice/Source/Joint/README.md}{the joint-family documentation};
their source conclusion is
\leanfile{S7_ParameterChoice/L3_Source.lean\#L14}{L3_Source.lean}.

Append four further child profiles and one root. Each new child calls
seven existing profiles. The root makes twelve calls: eight to existing
profiles and four to the new children. Exact route checks establish
the ordered partitions, child-domain inclusions, slope inequalities,
terminal inequalities, and positive output margins.
The root data are
\[
 p_{\mathrm{root}}=\frac12,\qquad
 I_{\mathrm{root}}=[1-10^{-10},1+10^{-10}],\qquad
 L_{\mathrm{root}}(\rho)=z:=\frac{625145177}{500000000}.
\]
\leannote{\texttt{rootDatum} and the route certificates,
\leanfile{S7_ParameterChoice/L4_Routes.lean\#L31}{L4_Routes.lean}}
The old profiles embed unchanged into the enlarged family. Their
proper-subhost implications remain valid with the additional induction
hypotheses, and the five new pieces satisfy the same interface.
The simultaneous induction of Section~\ref{sec:global-construction}
therefore supplies the root assertion. At ratio $1$, it gives a
constant $C\ge1$ such that every sufficiently large target $k$ and
every graph of density at least $1/2$ on at least $Ce^{zk}$ vertices
contains a monochromatic $K_k$.
See \texttt{dense\_root\_profile\_unconditional} in
\leanfile{S7_ParameterChoice/L5_Root.lean\#L14}{L5_Root.lean}.

\subsection{From the dense profile to the Ramsey number}

Take a complete graph on $N=\lceil(C+2)e^{zk}\rceil$ vertices.
One colour has density at least $1/2$, so the root assertion gives a
monochromatic $K_k$. Set $\gamma=z+10^{-9}=250058071/200000000$.
Since $\gamma>z$, the prefactor and integer rounding are absorbed
for all sufficiently large $k$:
\[
 R(k,k)\le\left\lceil(C+2)e^{zk}\right\rceil\le e^{\gamma k}.
\]
\leannote{\texttt{eventual\_exponent\_of\_dense\_profile},
\leanfile{S8_RamseyBound/L2_ComplementSymmetry.lean\#L23}{L2_ComplementSymmetry.lean}}
For the literal base, put $u=\gamma/2\in(0,1)$. The Taylor remainder
bound at order $14$ and exact rational arithmetic give
\[
 e^\gamma\le
 \left(\sum_{j=0}^{13}\frac{u^j}{j!}
       +\frac{15u^{14}}{14\cdot14!}\right)^2
 <\frac{69827131}{20000000}=3.49135655.
\]
\leannote{\texttt{exponent\_gap},
\leanfile{S8_RamseyBound/L3_BaseComparison.lean\#L20}{L3_BaseComparison.lean}}
Raising this comparison to the $k$th power proves
Theorem~\ref{thm:main}.

The linked development uses Lean~\cite{Lean} and
Mathlib~\cite{Mathlib}, both at version~4.35.0.
Its Ramsey definitions and foundational lemmas come from
RamseyLean~\cite{RamseyLean}, which formalizes~\cite{GNNW}.
We reuse the weighted-potential and profile-bank proof code and data
of~\cite{ZBS}, and the source-bound verification and analytic checking
library of~\cite{LW}; the latter includes interval-arithmetic code
originating in RamseyLean.
Exact inherited revisions and license notices are in the
\href{\leanrepo/NOTICE.md}{repository attribution record}.
The
\href{\leanrepo/Scripts/README.md}{reproduction instructions} describe
the source checks, project build, axiom audit, and built-in
\texttt{lake comparator} check against a separate finite-graph
specification. The comparison policy permits only
\texttt{propext}, \texttt{Classical.choice}, and \texttt{Quot.sound}.
